% ============================ 8. Conclusion =============================
\section{Conclusions}\label{sec:conclusion}

We set out to determine what computational structure and limitations
emerge when attention-like temporal weighting is compressed into a
minimal spiking dynamical unit, and to test, layer by layer, whether
temporal memory, attention-like weighting, routing and semantic binding
are distinct capabilities.  The frozen experiments support the following
conclusions, each with the evidence level stated:

\begin{itemize}
  \item \textbf{History dependence is established (Phase 1).}  The
    scalar membrane potential is not a sufficient state for TAN or for
    any window-carrying variant; the projection onto $h$ is not a
    lumpable reduction.  History dependence is a property of retained
    context, not of attention alone, and it is conditionally gated by
    prediction error.
  \item \textbf{Attention reorganizes local response geometry (Phase 2,
    supported with caveats).}  In the matched full representation, Full
    TAN shows higher effective rank and covariance scale in
    event-locked responses than the uniform-attention control, robustly
    across amplitude controls; this is a statement about local
    fluctuation geometry, not intrinsic dimensionality, and the shared
    observable frame shows no TAN-specific advantage over LIF.
  \item \textbf{Retrieval evidence is limited (Probe 1).}  Full TAN
    shows evidence of improved target-relevant retrieval under
    content-rich temporal context relative to the no-attention control,
    partially replicated across seeds, while a fixed delay line remains
    the strongest fixed-position memory baseline.
  \item \textbf{Query-conditioned routing is not implemented (Probe 2).}
    The scalar kernel fails the routing identifiability audit; changing
    the query while holding history fixed never changes which history
    tap is attended.  This is a mechanistic identifiability result, kept
    in full as \texttt{TERMINATED / AUDIT\_FAILED} in the archive.
\end{itemize}

The experiments therefore separate temporal memory, attention-like
saliency weighting, routing, and semantic binding as distinct
computational properties:
\begin{equation*}
  \boxed{\;
  \text{Memory} \;\neq\; \text{Weighting} \;\neq\;
  \text{Routing} \;\neq\; \text{Binding}
  \;},
\end{equation*}
and the mechanistic boundary of the current kernel can be stated
precisely: \emph{a minimal scalar temporal-attention mechanism can
generate history-dependent dynamics and dynamically reorganize
event-locked response structure, while remaining fundamentally limited
in its ability to perform query-conditioned content selection.}
Whether organizational extensions -- dendritic compartments,
opponent circuits, lateral or phase-based competition -- lift this
boundary is the legitimate next question (TAN-II); it is not answered,
and not claimed, by the present paper.
